\chapter{The Feed Handler}\label{ch:ll:the-feed-handler}

At the open the message rate jumps: on the simulator's busy instrument of this chapter, from about 570 messages a second at
midday to more than 9\,000 in the first seconds; on the consolidated feed of U.S. options, OPRA published a peak of 63.9
million messages in one second of July 2026 and 350\,000 in its busiest millisecond. A handler that falls behind lets its
socket buffer fill; the kernel drops what does not fit; line A loses a packet that line B has also lost; and the firm's book
is stale until the handler has recovered, by retransmission if the gap is small and recent, from the next snapshot if not.
This chapter builds the first component of Part~IV's trading system, the \omterm{def:ll:the-feed-handler:handler}{feed handler}: it decodes and arbitrates the two
lines of chapter~16, detects gaps and recovers from them, normalises every message into the firm's event record and publishes
it on the rings of chapter~12, and says at every moment whether the book it feeds can be trusted.

\section{Decode and arbitrate}

\begin{definition}[Feed handler]\label{def:ll:the-feed-handler:handler}
A \emph{feed handler}\index{feed handler} is the process that receives a venue's market-data feed (One Quant Book~1,
chapter~4) and turns it into the firm's internal events: it reads the redundant lines, arbitrates them into one sequenced
stream, detects and recovers lost messages, normalises each message into the firm's own record (Book~1, chapter~28), and
publishes the records to the components that build books and make decisions.
\end{definition}

\begin{omfigure}
\begin{tikzpicture}[font=\footnotesize,box/.style={draw,rounded corners=2pt,minimum height=0.7cm,align=center,inner sep=3pt},
  >=Stealth]
\node[box,fill=gray!10] (a) at (0,0.6) {line A};
\node[box,fill=gray!10] (b) at (0,-0.6) {line B};
\node[box,fill=omDef!20,minimum width=2.1cm] (arb) at (2.6,0) {arbitrate by\\ sequence};
\node[box,fill=omThm!15,minimum width=2.2cm] (gap) at (5.4,0) {gap? hold, wait\\ \qty{0.5}{\milli\second} for\\ the other line};
\node[box,fill=omDef!20,minimum width=2.1cm] (norm) at (8.3,0) {normalise\\ (48-byte event)};
\node[box,fill=omMeth!20] (ring) at (10.9,0) {firm.ring\\ (chapter~12)};
\node[box,fill=omThm!8] (retx) at (5.4,-2.0) {retransmission\\ server};
\node[box,fill=omThm!8] (snap) at (8.3,-2.0) {snapshot\\ channel};
\node[box,fill=gray!8,align=left] (flag) at (5.4,1.7) {stale-book flag, counters,\\ staleness intervals};
\draw[->] (a) -- (arb);
\draw[->] (b) -- (arb);
\draw[->] (arb) -- (gap);
\draw[->] (gap) -- (norm);
\draw[->] (norm) -- (ring);
\draw[->,omThm] (gap) -- node[left,font=\scriptsize]{request} (retx);
\draw[->,omThm,dashed] (retx.east) -- node[below,font=\scriptsize]{} ++(0.4,0) |- ([yshift=-0.2cm]gap.east);
\draw[->,omThm] (snap) -- node[right,font=\scriptsize]{if the server cannot} (norm);
\draw[gray,dashed] (gap) -- (flag);
\end{tikzpicture}

\omcaption{The build's \omterm{def:ll:the-feed-handler:handler}{feed handler}. Packets from both lines are merged by sequence number; a gap holds later messages and
waits briefly for the other line, then asks the retransmission server, then falls back to the next snapshot; every message
leaves as one fixed-size event on the firm's ring. While a gap is open, the \omterm{def:ll:the-feed-handler:stale}{stale-book flag} is up.}\label{fig:ll:the-feed-handler:design}
\end{omfigure}

The handler of the build is written first in Python, as the reference, then in C++20 and Rust, and all three publish exactly
the same events (compared by a 64-bit hash of their 48-byte records) on the same inputs. Its time is the arrival time recorded
with each packet, so that a recorded day replays exactly: the same inputs give the same events, the same staleness intervals
and the same counters. Arbitration is chapter~16's; what the handler adds is what to do while a sequence number is missing.

\section{Detect gaps and recover}

A missing sequence number is usually not lost, only late: it is on its way on the other line, or reordered by a switch.
The handler therefore holds the messages that arrive after it and waits a short time, here \qty{0.5}{\milli\second}, for the
missing ones. If they come, the held messages are released in order. If not, the handler asks the retransmission server for
the range (chapter~16; the simulator answers up to 1\,000 messages per request from its last 100\,000), and applies the answer
when it arrives, a round trip later. If the server cannot help, because the gap is older than its window or the server is
unreachable, the handler waits for the next snapshot of the instrument, loads it, discards the held messages it already
reflects and continues from the sequence number it gives, the procedure CME's snapshot messages support with their last
incremental sequence number processed.

\omcode{code/firm/feedhandler/cpp/firm_feedhandler.hpp}{213}{237}{Each packet: drop it if every message is old, deliver the
next expected message, hold the others; a held message opens a gap and starts its timer.}\label{lst:ll:the-feed-handler:incremental}

The snapshot path is the delicate one: the snapshot must be at least as recent as the book, the held messages it already
reflects must go, and a hole after it must be detected again.

\omcode{code/firm/feedhandler/firm_feedhandler.py}{190}{204}{Loading a snapshot in the Python reference.}\label{lst:ll:the-feed-handler:snapshot}

\begin{definition}[Stale-book flag, recovery latency]\label{def:ll:the-feed-handler:stale}
The \emph{stale-book flag}\index{stale-book flag} is the handler's public statement that the book built from its events may
differ from the venue's, raised when a gap is detected and lowered when it is closed. The \emph{recovery latency}\index{recovery
latency} of a gap is the time from its detection to the moment the book is correct again: the other line's delay, a
retransmission round trip, or the wait for and the loading of a snapshot.
\end{definition}

\Cref{fig:ll:the-feed-handler:minute} shows the handler at work on a simulated minute from the open, both lines impaired: 0.2\%
of packets lost independently on each line, bursts of loss, \omterm{def:ll:where-latency-comes-from:tail}{jitter} that reorders packets, and an outage on each line
overlapping for 15 milliseconds at 40 seconds. Out of 4\,232 gaps, 4\,224 are filled by the other line within the half
millisecond; the retransmission server closes the other eight within \qty{2.7}{\milli\second}. Without the server, seven gaps
need a snapshot, and with a snapshot every second the book stays stale for 0.3 to 1 second each time, five seconds of the
minute in all. Staleness is then not a property of the network but of the recovery path.

\begin{omfigure}
\begin{tikzpicture}
\begin{groupplot}[group style={group size=1 by 2,vertical sep=0.6cm,xlabels at=edge bottom},width=12cm,xmin=0,xmax=60,
  tick label style={font=\small},label style={font=\small},grid=major,grid style={gray!25},table/col sep=comma]
\nextgroupplot[height=3.6cm,ymin=0,ymax=1000,ylabel={msgs / 100 ms},xticklabels={}]
\addplot[omDef,const plot,fill=omDef!25,draw=omDef] table[x=t_s,y=messages]{figdata/low-latency/18-the-feed-handler/rate.csv} \closedcycle;
\nextgroupplot[height=3.2cm,ymin=0.4,ymax=2.6,ytick={1,2},yticklabels={{snapshot only},{with retrans.}},
  xlabel={seconds after the open},unbounded coords=jump]
\addplot[omMeth,line width=6pt,mark=|,mark size=4pt,mark options={line width=1.2pt}] table[x=t_s,y=retx]{figdata/low-latency/18-the-feed-handler/stale_bars.csv};
\addplot[omThm,line width=6pt] table[x=t_s,y=snapshot]{figdata/low-latency/18-the-feed-handler/stale_bars.csv};
\end{groupplot}
\end{tikzpicture}

\omcaption{A simulated minute from the open. Top: messages published per 100 ms, sixteen times busier at the open than at
midday. Bottom: the gaps that the other line could not fill, with the retransmission server (each closed within
\qty{3}{\milli\second}: a tick) and with snapshot recovery only (a snapshot a second: a bar from detection to recovery); the
4\,224 gaps filled by the other line are too short to draw. Data: \texttt{fig\_stale.py} (deterministic: the simulator and the handler
are).}\label{fig:ll:the-feed-handler:minute}
\end{omfigure}

\section{Normalise and publish}

Every message leaves the handler as one 48-byte event: its kind, side, instrument, sequence number, timestamp, order reference,
second reference (the new reference of a replace, the match number of an execution, the sequence number of a snapshot),
price and quantity. A normalised record has three virtues downstream. Its size is fixed, so it fits a ring slot
(chapter~12) and a consumer reads it without decoding. It is the same for every venue the firm reads, so the book builder
(chapter~19) and the strategies (chapter~20) are written once. And it carries the venue's sequence number, so a consumer can
tell exactly which events it has seen, and a snapshot's events can say which sequence number they reflect. The build's test
publishes a whole run on an SPSC ring of 64-byte slots and reads it back, hash for hash.

On the laptop the C++ handler takes about \qty{39}{\nano\second} per packet at the median, \qty{160}{\nano\second} at the p99
and \qty{290}{\nano\second} at the p99.9, measured over the minute of \cref{fig:ll:the-feed-handler:minute} with the clock read
around each input; the rare slow ones are gaps, whose held messages go through an ordered map, and retransmission answers.
At 9\,000 messages a second the handler is busy about 0.04\% of the time: the burst of this simulated instrument is nothing
for it. The real bursts are the ones of \cref{dat:ll:the-feed-handler:opra}.

\begin{dated}{2026-09}{How bursty is a feed?}\label{dat:ll:the-feed-handler:opra}
OPRA's published operating metrics for the consolidated U.S. options feed (August 2026 report) give, for July 2026, a peak of
63.9 million messages in one second, 7.9 million in 100 milliseconds, 0.89 million in 10 milliseconds and 0.35 million in
one millisecond: the busiest millisecond ran at 350 million a second, five and a half times the busiest second's rate. OPRA's capacity
notice says why subscribers plan on 10-millisecond peaks (they ``reflect system utilization during bursts of traffic''), and
that taking both redundant streams doubles the bandwidth.
\end{dated}

\section{Bursts, buffers and slow consumers}

\begin{definition}[Message burst, receive buffer overrun]\label{def:ll:the-feed-handler:burst}
A \emph{message burst}\index{message burst} is a period in which a feed's rate exceeds the rate its receiver sustains. A
\emph{receive buffer overrun}\index{receive buffer overrun} is the loss of datagrams that arrive while the socket's receive
buffer is full: the kernel drops them silently, and the handler sees only a gap.
\end{definition}

During a burst the difference between the arrival rate and the handler's service rate accumulates in the socket's receive
buffer, the M/M/1 queue of One Quant Book~4, chapter~8, with a hard wall: whatever does not fit is dropped. The buffer's
size is set with \texttt{SO\_RCVBUF}, and two details of Linux decide what it holds. The kernel ``doubles this value (to
allow space for bookkeeping overhead)'' and returns the doubled value (socket(7)), capped by \texttt{net.core.rmem\_max},
4~MiB on the laptop, so that asking for 4~MiB grants 8. And the buffer is charged with each datagram's whole memory footprint,
not its payload. \Cref{fig:ll:the-feed-handler:rcvbuf} measures what a stalled reader keeps: 50\,000 datagrams sent on
loopback while nobody reads, then the socket drained. Each 64-byte datagram costs 960 bytes of buffer and each 1\,000-byte
datagram 2\,305: the largest buffer the laptop grants holds 8\,738 small datagrams, a fortieth of the messages of OPRA's
busiest millisecond.

\begin{omfigure}
\begin{tikzpicture}
\begin{axis}[width=11.5cm,height=5.6cm,xmode=log,ymode=log,log basis x=2,xmin=20000,xmax=12000000,ymin=8,ymax=20000,
  xlabel={buffer granted by the kernel (bytes, log scale; twice the size asked)},ylabel={datagrams kept (log scale)},
  tick label style={font=\small},label style={font=\small},grid=major,grid style={gray!25},
  xtick={32768,131072,524288,2097152,8388608},xticklabels={32~KiB,128~KiB,512~KiB,2~MiB,8~MiB},
  legend image code/.code={\draw[#1] (0cm,0cm) -- (0.6cm,0cm);},
  legend columns=2,legend style={font=\footnotesize,at={(0.5,-0.3)},anchor=north,draw=none,fill=none,
  /tikz/every even column/.append style={column sep=8pt}},table/col sep=comma]
\addplot[omDef,very thick,mark=*,mark size=1.4pt] table[x=granted,y=held,restrict expr to domain={\thisrow{payload}}{0:100}]{figdata/low-latency/18-the-feed-handler/measured_rcvbuf.csv};
\addplot[omThm,very thick,mark=square*,mark size=1.4pt] table[x=granted,y=held,restrict expr to domain={\thisrow{payload}}{500:2000}]{figdata/low-latency/18-the-feed-handler/measured_rcvbuf.csv};
\legend{64-byte datagrams,1\,000-byte datagrams}
\end{axis}
\end{tikzpicture}

\omcaption{How many datagrams a UDP socket keeps for a reader that is not reading, by receive-buffer size: 50\,000 datagrams
sent on loopback during the stall. The kernel grants twice the size asked and charges each datagram its full memory
footprint. Measured on a laptop (Intel Core Ultra 7 155H) under WSL2 (kernel 6.18); \texttt{rmem\_max} 4~MiB. Data:
\texttt{bench\_feed.py}.}\label{fig:ll:the-feed-handler:rcvbuf}
\end{omfigure}

The handler is in turn a producer. Its consumers read the ring at their own pace, and a consumer slower than the feed (a
\omterm{def:ll:queues-ring-buffers-and-shared-memory:bp}{slow consumer}, chapter~12) must not slow the handler down: the broadcast ring never waits, and a lapped consumer
resynchronises. The two defences against bursts are therefore the same at both ends of the handler: enough buffer for the
burst the data says to expect, and a way to recover when even that is not enough.

\section{Monitoring a feed}

A handler that is wrong silently is worse than one that stops. The build counts packets, messages, duplicates, gaps, gaps
filled by the other line, retransmission requests and snapshot recoveries, and keeps every staleness interval with its start
and end. Four of these numbers deserve alarms: the rate of gaps on each line (a line going bad, while the other hides it), the
gaps filled by neither line (a common cause, chapter~16), any snapshot recovery (a book that was wrong for up to a snapshot
interval), and the receive queue's high-water mark (a burst closer to the buffer than expected). The \omterm{def:ll:the-feed-handler:stale}{stale-book flag} itself goes
to every consumer with each event, so that a strategy never quotes on a book that its handler knows to be wrong.

\section{Tutorial: the handler on impaired lines}\label{tut:ll:the-feed-handler:tut}

\begin{tutorial}
\textbf{Goal.} Run the handler in three languages on the same impaired lines, recover by retransmission and by snapshot, and
check the output against a clean run. \textbf{End state:} \cref{fig:ll:the-feed-handler:minute,fig:ll:the-feed-handler:rcvbuf},
the handling-time quantiles, and green tests in Python, C++ and Rust.
\begin{tutsteps}
\item \textbf{Make the fixture.} \texttt{make\_feed\_fixtures.py} runs Book~10's simulator for three seconds with both lines
impaired and writes the lines, the lossless stream (which the retransmission server holds) and the \omterm{def:ll:protocols-ii-binary-exchange-protocols:channels}{snapshot channel}.
\item \textbf{Three runs, three languages.} Clean, impaired with the server, impaired without it: the Python reference writes
the expected counts, hashes and staleness, and the C++ and Rust tests reproduce them exactly.
\item \textbf{Check what recovery means.} With retransmission, the impaired run publishes exactly the clean run's events. With
a snapshot, it publishes the snapshot instead of the lost messages, then exactly the clean run's events after it, and the book
at the end of the run is the same.
\item \textbf{Measure} with \texttt{python bench\_feed.py}: the minute of \cref{fig:ll:the-feed-handler:minute} through the
C++ handler with the clock read around each input (the script checks that it publishes what the reference publishes), and the
stalled-reader experiment of \cref{fig:ll:the-feed-handler:rcvbuf}.
\end{tutsteps}
\textbf{What to change next.} Shorten the wait for the other line to \qty{50}{\micro\second} and count the retransmission requests
that follow; lengthen the snapshot interval to five seconds and read the staleness again.
\end{tutorial}

\section{Build: the feed handler}\label{bld:ll:the-feed-handler:feedhandler}

\begin{build}
\bfield{Purpose} The first stage of the trading system: from the venue's lines to the firm's events. The book builder of
chapter~19 consumes its events, the strategy engine of chapter~20 its \omterm{def:ll:the-feed-handler:stale}{stale-book flag}, and the monitoring of chapter~24 its
counters.
\bfield{Interface} Python reference \texttt{firm\_feedhandler}: \texttt{Handler(retx, gap\_timeout\_ns, rtt\_ns).run(events)},
\texttt{.events}, \texttt{.hash}, \texttt{.stale}, \texttt{.counters}; \texttt{RetxServer}, \texttt{merge}, \texttt{recorded},
\texttt{normalise}, \texttt{book}. C++20 \texttt{firm::feed2::Handler(retx, timeout, rtt)} with \texttt{run(a, b, snapshots)},
\texttt{on\_event}, \texttt{hash}, \texttt{stale}, counters and optional timing hooks; \texttt{RetxServer}; \texttt{recorded}.
Rust \texttt{firm\_feedhandler::Handler} with the same.
\bfield{Rules} Each sequence number is published once, in order; held messages are released only in order; a gap waits for the
other line, then retransmission (one request per hole, answers capped by the server), then a snapshot; the handler is a
function of its timestamped inputs; the \omterm{def:ll:the-feed-handler:stale}{stale-book flag} is up exactly while a gap is open; no allocation per packet without a
gap.
\bfield{Acceptance tests} \texttt{code/firm/feedhandler/}: on the shared fixture, clean, retransmission and snapshot runs with
identical event counts, hashes, staleness intervals and counters in Python, C++ and Rust; the retransmission run equal to the
clean run event for event; the snapshot run equal to the clean run after the snapshot and in its final book; publication on
\texttt{firm.ring} read back hash for hash; arbitration on scripted packets.
\bfield{Stretch} One request for all the holes of a gap; per-instrument recovery on a multi-instrument feed; reading the lines
from sockets (chapter~13's \omterm{def:ll:linux-tuning:polling}{busy polling}) with the receive queue's high-water mark as a counter.
\end{build}

\begin{omsources}
\item Linux man page \texttt{socket(7)} (\texttt{SO\_RCVBUF}); Linux kernel documentation, ``Documentation for
/proc/sys/net/'' (\texttt{rmem\_max}, \texttt{rmem\_default}, \texttt{netdev\_max\_backlog}).
\item OPRA, \emph{Key Operating Metrics of U.S. Options Securities Information Processor} (August 2026); SIAC, \emph{Revised
OPRA Capacity Projections} (September 2025).
\item CME Group Client Systems Wiki, MDP~3.0: ``Market Data Snapshot -- Full Recovery''.
\item The simulator's protocols and services, \path{code/firm/exchsim/PROTOCOL.md} (One Quant Book~10, chapter~26).
\end{omsources}

\section{Exercises}

\begin{exercise}[$\star$]\label{exo:ll:the-feed-handler:1}
The handler expects sequence number 1\,000 and receives packets starting at 1\,003 (3 messages) on line A, then 1\,000 (3
messages) on line B. What does it publish, in what order, and what is the staleness interval if the second packet arrives
\qty{40}{\micro\second} after the first?
\end{exercise}

\begin{exercise}[$\star$]\label{exo:ll:the-feed-handler:2}
A handler asks for a 256~KiB receive buffer. What does \texttt{getsockopt} return, and how many 64-byte datagrams will the
buffer hold, by the measurement of \cref{fig:ll:the-feed-handler:rcvbuf}?
\end{exercise}

\begin{exercise}[$\star$]\label{exo:ll:the-feed-handler:3}
With snapshots every 500 milliseconds and a snapshot that takes 5 milliseconds to receive, what are the mean and worst recovery
latencies of an unrecoverable gap?
\end{exercise}

\begin{exercise}[$\star\star$]\label{exo:ll:the-feed-handler:4}
Why does the handler wait for the other line before asking for a retransmission? What does a shorter wait cost, and a longer
one?
\end{exercise}

\begin{exercise}[$\star\star$]\label{exo:ll:the-feed-handler:5}
A burst arrives at 2 million datagrams a second for 30 milliseconds; the handler serves 1.2 million a second. How many
datagrams queue, and what receive buffer (asked) holds them if each is charged 960 bytes?
\end{exercise}

\begin{exercise}[$\star\star$]\label{exo:ll:the-feed-handler:6}
After a snapshot recovery the handler's event stream differs from the clean run's. In what sense is the output nevertheless
``identical once recovered'', and how does the build test it?
\end{exercise}

\begin{exercise}[$\star\star\star$]\label{exo:ll:the-feed-handler:7}
\emph{Coding.} Make the handler ask for all the holes of a gap in one request (a list of ranges), in the Python reference and
the C++ handler, and check on the minute of \cref{fig:ll:the-feed-handler:minute} that the worst retransmission staleness falls.
\end{exercise}

\begin{exercise}[$\star\star\star$]\label{exo:ll:the-feed-handler:8}
\emph{Find the flaw.} ``When our handler detects a gap it drops everything it holds and waits for the next snapshot; that way
the book is always consistent.''
\end{exercise}

\section{Problem: The Opening Burst}

\begin{problem}[{Weekend problem --- buffers for the open, staleness when they fail}]\label{pb:ll:the-feed-handler:1}
A handler reads one line of a feed whose opening burst reaches 1.5 million datagrams a second for 20 milliseconds, one
message per datagram; the handler serves one datagram in \qty{2}{\micro\second} (decoding, book update, publication). Take the
kernel's charge per 64-byte datagram and its doubling of \texttt{SO\_RCVBUF} from \cref{fig:ll:the-feed-handler:rcvbuf}, a
maximum \texttt{rmem\_max} of 4~MiB, and a snapshot every second.

\medskip\noindent\textbf{Part I --- The queue.}
\begin{enumerate}
\item What is the handler's service rate, and its utilisation during the burst?
\item How many datagrams are queued at the end of the burst?
\item How long after the burst does the queue take to drain if the rate falls to 200\,000 a second?
\item What is the queueing delay of the last datagram of the burst?
\end{enumerate}

\medskip\noindent\textbf{Part II --- The buffer.}
\begin{enumerate}[resume]
\item How many bytes of receive buffer must the kernel grant to hold the queue?
\item What must be asked with \texttt{SO\_RCVBUF}, and is it allowed on this machine?
\item With the largest buffer allowed, how many datagrams are lost?
\item What else could hold the burst?
\end{enumerate}

\medskip\noindent\textbf{Part III --- When it fails.}
\begin{enumerate}[resume]
\item Can the retransmission server fill the loss?
\item If not, what staleness follows, on average and at worst?
\item During that staleness, how many datagrams arrive and must be held?
\item What should the strategy do meanwhile?
\end{enumerate}

\medskip\noindent\textbf{Part IV --- The verdict.}
\begin{enumerate}[resume]
\item State the \emph{named result}: the receive buffer that absorbs the burst without loss, and the staleness a snapshot
interval implies when it does lose.
\item How much faster would the handler have to be to need no queue at all?
\item Why does reading both lines double the problem?
\item What does OPRA's millisecond peak say about a single-threaded handler for a whole options feed?
\item What would you measure in production to size the buffer?
\item Which counter tells you that a burst came close to the limit?
\item Why is a snapshot interval a business decision of the venue with a latency cost for the firm?
\item In one sentence: what does a \omterm{def:ll:the-feed-handler:handler}{feed handler} do when it cannot keep up?
\end{enumerate}
\end{problem}

\section{Interview questions}

\begin{interviewq}[$\star$ \iqroles{developer}]\label{iq:ll:the-feed-handler:1}
What does a market-data \omterm{def:ll:the-feed-handler:handler}{feed handler} do, from the network card to the strategy?
\end{interviewq}

\begin{interviewq}[$\star\star$ \iqroles{developer}]\label{iq:ll:the-feed-handler:2}
A gap appears on line A. Walk through what the handler does until the book is correct again.
\end{interviewq}

\begin{interviewq}[$\star\star$ \iqroles{developer}]\label{iq:ll:the-feed-handler:3}
How would you size a UDP socket's receive buffer for a feed? What does the kernel do with the value you set?
\end{interviewq}

\begin{interviewq}[$\star\star$ \iqroles{developer}]\label{iq:ll:the-feed-handler:4}
How do you know, downstream, that the book you are reading is stale?
\end{interviewq}

\begin{interviewq}[$\star\star$ \iqroles{developer}]\label{iq:ll:the-feed-handler:5}
Why normalise every venue's messages into one internal record? What does it cost?
\end{interviewq}

\begin{interviewq}[$\star\star\star$ \iqroles{developer, researcher}]\label{iq:ll:the-feed-handler:6}
How would you test a \omterm{def:ll:the-feed-handler:handler}{feed handler}'s recovery paths before production, and prove that it recovers to the right book?
\end{interviewq}
